\begin{resumo}

O volume de dados coletado cresce em ritmo acelerado, e com ele o uso de aprendizado de
máquina para extrair informação desses registros. Em diversas aplicações, o resultado
procurado é a proporção de cada classe dentro de um grupo, como a fração de infectados
em uma população ou a de menções negativas sobre um produto. Classificar cada elemento e
contar os rótulos resolve o problema enquanto os dados de treinamento e os de aplicação
vierem da mesma distribuição, hipótese que a prática viola com frequência. Quando a
proporção se desloca, o classificador erra sempre na mesma direção e a contagem fica
enviesada. Dessa limitação nasce a quantificação, cujos métodos estimam a proporção de
forma direta. As implementações disponíveis em Python estão dispersas em pacotes com
formatos próprios, o que obriga o usuário a adaptar dados e métodos a cada biblioteca e
dificulta a comparação entre elas. Neste trabalho, é apresentada a \texttt{mlquantify},
uma biblioteca em que 45 métodos de quantificação são reunidos sob a interface do
\textit{scikit-learn}. Na arquitetura, são combinados \textit{mixins} e \textit{tags}
para acomodar exigências distintas, a decomposição multiclasse é aplicada de forma
automática e novos métodos podem ser compostos a partir de representação, função de
perda e solucionador. São disponibilizados ainda protocolos de avaliação, métricas de
erro e busca de hiperparâmetros guiada pelo erro de quantificação. A biblioteca foi
comparada com a \texttt{QuaPy}, a \texttt{QuantificationLib} e a \texttt{qunfold} em 28
métodos e 132 conjuntos de dados, sempre sobre as mesmas amostras. O erro foi avaliado
por teste de equivalência, com margem de um ponto percentual de prevalência, e o custo
pela razão de tempo e de memória entre as implementações. Em 38 dos 39 pares de método e
biblioteca o intervalo de confiança da diferença de erro coube dentro da margem, e a
única exceção favorece a biblioteca proposta. A razão mediana de custo ficou em 0,967 no
tempo e em 0,816 na memória. Entre as bibliotecas comparadas, a proposta é a de
cobertura mais ampla de métodos, com erro equivalente e custo da mesma ordem de
grandeza, sob uma interface única.

\vspace{\onelineskip}
\noindent
\textbf{Palavras-chaves}: Quantificação. Aprendizado de máquina. Biblioteca de software. Estimativa de prevalência. Avaliação experimental.
\end{resumo}
